% TODO
% Response to referee 3, comment 1
\begin{enumerate}
\item We can just answer honestly that there is numerical evidence for the material being proximately quantum critical but no experimental evidence as it is a pretty new/understudied material.
\item In some sense one could argue that the spectrum of collective excitations is a kind of probe of quantum criticality, so that our own observations provide physical evidence.
\item It is essential to note that no conclusions of our work related to assigining the mode we observe to a Higgs mode relies on the material being quantum critical.
The quantum criticality is a reasonable, but untested hypothesis intended to motivate future work on this and related materials.
\textcolor{red}{Is this strategically sound. also, probably good to clarify our interpretation does not ``rest" on criticality. We also point out quasi-1d physics.}
\end{enumerate}


We thank the referee for their close consideration of our paper and very helpful comments to improve our manuscript.

The primary evidence for \ce{CuBr2} being quantum critical is numerical.
\citet{furukawa_ground-state_2012} presents a comprehensive study of the spin-$1/2$ frustrated XXZ chain, which is a good description of the low-energy magnetic Hamiltonian of \ce{CuBr2} due to the quasi-1d arrangement of the \ce{CuBr4} spin planes.
Figure $1$ in that reference describes the numerical phase diagram as a function of the ratio of the magnetic exchange parameters $J_1/J_2$ and the anisotropy gap $\Delta$. 
From \textcolor{red}{[Ajesh: what papers are these numbers from?]}, \ce{CuBr2} lies near \textcolor{red}{$\Delta=\qty{?}$, $J_1/J_2=\qty{?}$} in that parameter space, which is close to the quantum critical point separating the observed vector chiral phase from a paramagnetic Haldane dimer phase.

Importantly, however, we would like to stress that our interpretation does not ``rest'' on \ce{CuBr2} being quantum critical.
Certainly the identification of the slow \gls{shg} oscillation as a coherent Higgs excitation is independent of the system being quantum critical.
The fact that \ce{CuBr2} is located close to a quantum critical point was invoked both to provide a plausible explanation for the mechanism stabilizing the mode in this material, as well as to motivate future studies into the many aspects of \ce{CuBr2} that cannot be understood entirely in a single paper.
Given the strong one-dimensionality of \ce{CuBr2}, as well as the complex phase diagram of the XXZ spin-$1/2$ model to which \ce{CuBr2} belongs, we believe this invocation is reasonable.

\begin{changesmade}
Added a sentence to the last paragraph of the discussion section: ``Further studies regarding the quantum critical nature of \ce{CuBr2} and related multiferroic compounds will be of great interest.''
\end{changesmade}

\textcolor{red}{Revisit: does Nuh think this is strategically OK? Also ask Ajesh for numbers}
